\PassOptionsToPackage{quiet}{xeCJK}
\documentclass[answer]{THUExam}
% \documentclass[]{THUExam}
\usepackage{HTNotes-math}
\usepackage{pifont}
\usepackage{ulem}

\usepackage{tikz}
\usetikzlibrary{positioning, calc, shapes.geometric, shapes, shapes.multipart, arrows.meta, arrows, decorations.markings, external, trees}

% =================================================
%       PDF信息
% =================================================
% PDF信息里的标题栏
\title{高等数学A (二)}
% PDF信息里的作者栏
\author{}
% PDF信息里的主题
\Subject{作业}
% PDF信息里的关键词
\Keywords{高等数学}

% =================================================
%       试卷头信息
% =================================================
% 试卷头里的年份
\Year{2021--2022}
% 试卷头里的学期
\Semester{秋季}
% 试卷头里的课程
\Course{高等数学A}
% 考试时间
\Time{}
% 试卷头里的类型，如A/B/模拟等
\Type{A}

\begin{document}
\vspace{1cm}
% 生成试卷表头
\makehead

\vspace{1cm}

\makepart{填空题} {每小题错误扣~$1$~分}

\begin{problem}
试说出下列各微分方程的阶数:
\begin{enumerate}
    \item $x\left(y^{\prime}\right)^2-2 y y^{\prime}+x=0$, \quad\fillin{$\quad 1\quad $}阶;
    \item $x^2 y^{\prime \prime}-x y^{\prime}+y=0$, \quad\fillin{$\quad 2\quad $}阶;
    \item $x y^{\prime \prime \prime}+2 y^{\prime \prime}+x^2 y=0$, \quad\fillin{$\quad 3\quad $}阶;
    \item $(7 x-6 y) \mathrm{d} x+(x+y) \mathrm{d} y=0$, \quad\fillin{$\quad 1\quad $}阶;
    \item $L \frac{\mathrm{d}^2 Q}{\mathrm{~d} t^2}+R \frac{\mathrm{d} Q}{\mathrm{~d} t}+\frac{Q}{C}=0$, \quad\fillin{$\quad 2\quad $}阶;
    \item $\frac{\mathrm{d} \rho}{\mathrm{d} \theta}+\rho=\sin ^2 \theta$, \quad\fillin{$\quad 1\quad $}阶.
\end{enumerate}
\end{problem}
\bigskip
\begin{note}
    参照 《高等数学》(上册), Page 291, 关于微分方程阶的定义:
    \begin{quote}
        微分方程中所出现的未知函数的最高阶导数的阶数, 叫做微分方程的阶.
    \end{quote}
\end{note}
\bigskip
\newpage
\makepart{解答题}{每小题错误扣~$2$~分}

\begin{problem}
指出下列各题中的函数是不是所给微分方程的解:
\begin{enumerate}
    \item $y^{\prime \prime}-\left(\lambda_1+\lambda_2\right) y^{\prime}+\lambda_1 \lambda_2 y=0,\quad y=C_1 e^{\lambda_1 x}+C_2 e^{\lambda_2 x}$.
\end{enumerate}
\end{problem}
\bigskip
\begin{note}
    参照 《高等数学》(上册), Page 291, 关于微分方程解的定义:
    \begin{quote}
        ...然后找出满足微分方程的函数 (解微分方程), 就是说, 找出这样的函数, 把这函数代人微分方程能使该方程成为恒等式. 这个函数就叫做该微分方程的解.
    \end{quote}
\end{note}
\begin{solution}
    \begin{enumerate}
        \item 分别计算 $y$, $y'$ 及 $y''$;
              \begin{enumerate}
                  \item $y=C_1 e^{\lambda_1 x}+C_2 e^{\lambda_2 x}$;
                  \item $y'=\lambda_1C_1  e^{\lambda_1 x}+\lambda_2 C_2 e^{\lambda_2 x}$;
                  \item $y''= \lambda_1^2 C_1e^{\lambda_1 x}+\lambda_2^2 C_2 e^{\lambda_2 x}$;
              \end{enumerate}
        \item 将$y$, $y'$ 及 $y''$ 代入方程;
              $$
                  \begin{aligned}
                        & \,\,y^{\prime \prime}-\left(\lambda_1+\lambda_2\right) y^{\prime}+\lambda_1 \lambda_2 y                                                                                                                                          \\
                      = & \,\, \uline{\lambda_1^2 C_1 e^{\lambda_1 x}}+\uwave{\lambda_2^2 C_2 e^{\lambda_2 x}}-\uline{\lambda_1\left(\lambda_1+\lambda_2\right) C_1 e^{\lambda_1 x}}-\uwave{\lambda_2\left(\lambda_1+\lambda_2\right) C_2 e^{\lambda_2 x}} \\
                        & \,\,+ \uline{\lambda_1 \lambda_2 C_1 e^{\lambda_1 x}}+\uwave{\lambda_1 \lambda_2 C_2 e^{\lambda_2 x}},                                                                                                                           \\
                      = & \,\, \uline{C_1 \left(\lambda_1^2 -\lambda_1\left(\lambda_1+\lambda_2\right) + \lambda_1 \lambda_2 \right)  e^{\lambda_1 x}}                                                                                                     \\
                        & \,\,+\uwave{C_2 \left(\lambda_2^2 -\lambda_2\left(\lambda_1+\lambda_2\right) + \lambda_1 \lambda_2 \right)  e^{\lambda_2 x}}                                                                                                     \\
                      = & \,\,0+0=\,\,0.
                  \end{aligned}
              $$
    \end{enumerate}
    因此该函数是所给微分方程的解.
\end{solution}
\bigskip

\begin{problem}
在下列各题中, 验证所给二元方程所确定的函数为所给微分方程的解:
\begin{enumerate}
    \item $(x-2 y) y^{\prime}=2 x-y, x^2-x y+y^2=C$;
    \item $(x y-x) y^{\prime \prime}+x y^{\prime 2}+y y^{\prime}-2 y^{\prime}=0, y=\ln (x y)$.
\end{enumerate}
\end{problem}
\begin{note}
    参照 , 关于微分方程解的定义:
    \begin{quote}
        ...然后找出满足微分方程的函数 (解微分方程), 就是说, 找出这样的函数, 把这函数代人微分方程能使该方程成为恒等式. 这个函数就叫做该微分方程的解.
    \end{quote}
    并参照《高等数学》(上册), Page 98-101, 隐函数求导法则.
\end{note}
\begin{solution}
    \begin{enumerate}
        \item  在方程 $x^2-x y+y^2=C$ 两端对 $x$ 求导, 得
              $$
                  2 x-\left(y+x \frac{\mathrm{d}y}{\mathrm{d}x}\right)+2 y \frac{\mathrm{d}y}{\mathrm{d}x}=0,
              $$

              即 $(x-2 y) y^{\prime}=2 x-y$. 所以给出的二元方程所确定的函数是微分方程的解.
        \item 在方程 $y=\ln (x y)$ 两端对 $x$ 求导, 得
              $$
                  \frac{\mathrm{d}y}{\mathrm{d}x}=\frac{y+x \frac{\mathrm{d}y}{\mathrm{d}x}}{x y},
              $$

              即 $(x y-x) \frac{\mathrm{d}y}{\mathrm{d}x}-y=0$, 再在上式两端对 $x$ 求导, 得
              $$
                  \left(y+x \frac{\mathrm{d}y}{\mathrm{d}x}-1\right) \frac{\mathrm{d}y}{\mathrm{d}x}+(x y-x) \frac{\mathrm{d}^2y}{\mathrm{d}x^2}-\frac{\mathrm{d}y}{\mathrm{d}x}=0,
              $$

              即 $(x y-x) y^{\prime \prime}+x y^{\prime 2}+y y^{\prime}-2 y^{\prime}=0$. 所以给出的二元方程所确定的函数是所给微分方程的解.
    \end{enumerate}
\end{solution}

\begin{problem}
求下列微分方程的通解:
\begin{itemize}
    \item (2) $3 x^2+5 x-5 y^{\prime}=0$;
    \item (6) $\frac{\mathrm{d} y}{\mathrm{~d} x}=10^{x+y}$;
    \item (8) $\cos x \sin y \mathrm{~d} x+\sin x \cos y \mathrm{~d} y=0$;
\end{itemize}
\end{problem}
\begin{note}
    参照《高等数学》(上册), Page 295, 可分离变量的微分方程的分析方法, 也需要参照Page 180-207 得积分表, 第一, 第二类换元积分法和分部积分法.
\end{note}
\begin{solution}
    \begin{itemize}
        \item (2) 这里可以移项为 $5 y^{\prime}=3 x^2+5 x$ 后直接不定积分求解:
              $$
                  5 y=x^3+\frac{5}{2} x^2+C.
              $$
        \item (6) 进行分离变量,  $10^{-y} \mathrm{~d} y=10^x \mathrm{~d} x$, 使用可分离变量的偏微分方程的分析方法得
              $$
                  -\frac{10^{-y}}{\ln 10}=\frac{10^x}{\ln 10}+C,
              $$
              可简化为
              $$
                  10^x+10^{-y}=C.
              $$
        \item (8) 进行分离变量, 得 $\frac{\cos y}{\sin y} \mathrm{~d} y=-\frac{\cos x}{\sin x} \mathrm{~d} x$,两端积分得
              $$
                  \ln |\sin y|=-\ln |\sin x|+\ln C_1,
              $$
              即
              $$
                  \sin y\sin x=C.
              $$
    \end{itemize}
\end{solution}


\begin{problem}
求下列微分方程满足所给初值条件的特解:
\begin{itemize}
    \item (1) $y^{\prime}=\mathrm{e}^{2 x-y},\left.y\right|_{x=0}=0$;
    \item (4) $\cos y \mathrm{~d} x+\left(1+\mathrm{e}^{-x}\right) \sin y \mathrm{~d} y=0,\left.y\right|_{x=0}=\frac{\pi}{4}$;
    \item (6) $\frac{1}{\rho} \frac{\mathrm{d} \rho}{\mathrm{d} \theta}+\frac{\rho^2+1}{\rho^2-1} \cot \theta=0,\left.\rho\right|_{\theta=\frac{\pi}{6}}=3$;
\end{itemize}
\end{problem}
\begin{note}
    参照《高等数学》(上册), Page 295, 可分离变量的微分方程的分析方法, 也需要参照 Page 292, 特解和初值问题的定义.
\end{note}

\begin{solution}
    \begin{itemize}
        \item (1) 分离变量, 得 $\mathrm{e}^y \mathrm{~d} y=\mathrm{e}^{2 x} \mathrm{~d} x$, 两端积分得
              $$
                  \mathrm{e}^y=\frac{1}{2} \mathrm{e}^{2 x}+C,
              $$

              由 $\left.y\right|_{x=0}=0$, 得 $1=\mathrm{e}^0=\frac{1}{2} \mathrm{e}^0+C$, 故 $C=\frac{1}{2}$, 即得 $\mathrm{e}^y=\frac{1}{2}\left(\mathrm{e}^{2 x}+1\right)$, 于是所求特解为 $y=\ln \frac{\mathrm{e}^{2 x}+1}{2}$.
        \item (4) 分离变量,得 $\frac{\mathrm{e}^x}{\mathrm{e}^x+1} \mathrm{~d} x=-\tan y \mathrm{~d} y$,两端积分得
              $$
                  \ln \left(\mathrm{e}^x+1\right)=\ln |\cos y|+\ln C_1,
              $$

              即 $\mathrm{e}^x+1=C \cos y$. 代人初值条件: $x=0, y=\frac{\pi}{4}$, 有 $2=C \cdot \frac{\sqrt{2}}{2}$, 得 $C=2 \sqrt{2}$, 于是
              $$
                  \mathrm{e}^x+1=2 \sqrt{2} \cos y,
              $$

              即 $\left(\mathrm{e}^x+1\right) \sec y=2 \sqrt{2}$ 为所求特解.
        \item (6) 分离变量得到 $\frac{1}{\rho} \cdot \frac{\rho^2-1}{\rho^2+1} \mathrm{~d} \rho=-\cot \theta \mathrm{d} \theta$, 等式两端积分, 对于左端, 注意使用第一类换元积分法, 得到
              $$
                  \begin{aligned}
                      \text{LHS} & = \int \frac{\rho}{\rho^2} \cdot \frac{\rho^2-1}{\rho^2+1} \mathrm{~d} \rho                                                                                         \\
                                 & =\frac{1}{2} \int \frac{\rho^2-1}{\rho^2\left(\rho^2+1\right)} \mathrm{d} \rho^2=\frac{1}{2} \int\left(\frac{2}{\rho^2+1}-\frac{1}{\rho^2}\right) \mathrm{d} \rho^2 \\
                                 & =\ln \left(\rho^2+1\right)-\frac{1}{2} \ln \rho^2+\ln C_1=\ln \left(\rho+\frac{1}{\rho}\right)+\ln C_1,
                  \end{aligned}
              $$
              对于右端, 也是利用第一类换元积分法, 得到
              $$
                  \text{RHS}=\int(-\cot \theta) \mathrm{d} \theta=\int\left(-\frac{\cos \theta}{\sin \theta}\right) \mathrm{d} \theta=-\ln \sin \theta+\ln C_2,
              $$
              也就是
              $$
                  \rho+\frac{1}{\rho}=\frac{C}{\sin \theta} .
              $$

              代人初值条件 $\theta=\frac{\pi}{6}, \rho=3$, 得 $C=\frac{5}{3}$, 故所求特解为 $\rho+\frac{1}{\rho}=\frac{5}{3 \sin \theta}$.
    \end{itemize}
\end{solution}

\begin{problem}
一曲线通过点 $(2,3)$, 它在两坐标轴间的任一切线线段均被切点所平分, 求这曲线方程.
\end{problem}
\begin{solution}
    设曲线方程为 $y=y(x)$, 切点为 $(x, y)$. 依条件, 切线在 $x$ 轴与 $y$ 轴上的交点分别为 $(2 x,0)$ 与 $(0,2 y)$. 于是, 切线的斜率
    $$
        y^{\prime}=\frac{2 y-0}{0-2 x}=-\frac{y}{x} .
    $$
    由于曲线过点 $(2,3)$, 所以曲线满足条件
    $$
        y(2)=3.
    $$
    因此, 可以对曲线方程建立微分方程为
    $$
        \left\{
        \begin{aligned}
            y^{\prime} & =-\frac{y}{x}, \\
            y(2)       & =3.
        \end{aligned}
        \right.
    $$
    对该微分方程进行分离变量得
    $$
        \frac{\mathrm{d} y}{y}=-\frac{\mathrm{d} x}{x} .
    $$

    积分得 $\ln |y|=-\ln |x|+\ln C_1$, 即 $x y=C$. 代人初值条件 $x=2, y=3$, 得 $C=6$, 故曲线方程为 $x y=6$.
\end{solution}
\end{document}
